Load forecasts for battery scheduling are usually chosen by accuracy metrics such as the mean absolute error. We ask whether the most accurate forecast also saves the most money. In a pre-registered study of \res{nKept} non-residential buildings from the Building Data Genome 2 data set, a behind-the-meter battery was scheduled day-ahead from five standard forecasters under a time-of-use tariff with and without a monthly demand charge, and every schedule was costed on the buildings' actual 2017 load. Under the demand charge, the ranking of forecasters by accuracy did not agree with their ranking by saving (mean Kendall $\tau=\res{honeMean}$ with open-loop execution and $\res{honebMean}$ with a real-time peak guard), and the most accurate forecaster was the most valuable in only \res{mostAccOpen}--\res{mostAccGuard}\% of building-months, below the 20\% expected by chance. The least accurate model on average, a 0.9-quantile forecaster, saved the most. At equal error, a one-hour timing error destroyed the saving while a uniform over-forecast of about 10\% kept most of it. How the schedule was executed mattered more than the forecaster: executed open loop, the schedules of the four point forecasters raised the bill above the no-battery bill in \res{negOpenFourLo}--\res{negOpenFourHi}\% of buildings (\res{negT1F5}\% for the quantile forecaster), while a simple guard turned all of them into savings and beat a forecast-free reactive rule in \res{hfourbBeatPct}\% of buildings. Without a demand charge, all forecasters captured the full saving. Forecasts for storage under demand charges should be judged by the bill they produce, not by accuracy.
